\documentclass[11pt,a4paper]{article}
\usepackage{curso}
\begin{document}
\portada{Evaluación --- Unidad}{11}{Quiz: software, agentes y fronteras}
\noindent Preguntas guía: ¿pueden los agentes acelerar la coevolución de algoritmos y hardware? ¿Qué te convencería de que una ventaja termodinámica es real?
\section*{Preguntas}
\begin{enumerate}
\item En THRML, los nodos de un mismo bloque deben: (a) interactuar todos (b) no interactuar entre sí (c) ser continuos (d) estar fijados.
\item Una compuerta PNOT$(p)$ con $p=\tfrac12$ produce: (a) la identidad (b) un NOT determinista (c) la uniforme (d) un estado fijo.
\item La identidad de desplazamiento logarítmico expresa $\partial_\theta G$ como: (a) una integral (b) la diferencia de dos compuertas ejecutables (c) un hessiano (d) cero.
\item Componer PIsing por aristas en un anillo: (a) es Gibbs exacto (b) no da la distribución de Boltzmann del anillo (c) diverge (d) es determinista.
\item Un termalizador realiza un núcleo: (a) fijando entradas, termalizando y leyendo salidas (b) con una red densa (c) por enumeración (d) en una GPU.
\item Con núcleos que contraen, el error de compilación con la profundidad: (a) crece exponencialmente (b) crece linealmente siempre (c) se satura (d) es nulo.
\item En GRPO la ventaja de un intento es: (a) su recompensa (b) recompensa menos media del grupo, sobre la desviación (c) el valor de una red crítica (d) la KL.
\item Thermo-RSI demostró: (a) descubrimiento de algoritmos nuevos (b) mejora al reproducir experimentos clásicos (c) ejecución en Z1 (d) automejora recursiva.
\item (Problema) Recompensa de una solución que corre, no reproduce la métrica y cumple la rúbrica ($r=0{,}7\,\text{exec}+0{,}3\,\text{rúbrica}$).
\item (Problema) Giros/ns implicados por 269\,568 pbits a 50 MHz.
\item (Problema) Cota del error estacionario con $\rho_0=0{,}85$ y error por paso $0{,}0155$.
\item (Problema) Ventajas de GRPO para recompensas $(1;\,0{,}4;\,0;\,0{,}7)$.
\end{enumerate}
\newpage
\section*{Clave}
1 (b) · 2 (c) · 3 (b) · 4 (b) · 5 (a) · 6 (c) · 7 (b) · 8 (b).
\begin{enumerate}\setcounter{enumi}{8}
\item $0{,}7\cdot0{,}15+0{,}3\cdot1=0{,}405$.
\item $\approx13\,478$.
\item $0{,}0155/0{,}15\approx0{,}10$.
\item $(+1{,}28;\,-0{,}34;\,-1{,}42;\,+0{,}47)$.
\end{enumerate}
Aprobación 8/12; fallos en 4, 6 u 8 $\Rightarrow$ repasar 11.2, 11.3, 11.5.
\end{document}
